\documentclass[12pt]{article}
\usepackage[utf8]{inputenc}
\usepackage[T1]{fontenc}
\usepackage{hyperref}
\title{Yiningayi Language Resources for K-12 Curriculum Development}
\author{Compiled by Claude Sonnet 5.5 (claude-sonnet-5-5) for Andr\'as Kornai}
\date{}
\begin{document}
\maketitle

\section{Introduction}
Yiningayi (ISO 639-3 \texttt{ygi}; also Iningai, Mootaburra, Muttaburra, Yiningay) is an extinct Australian language of Queensland. The AIATSIS record places it west of the Dividing Range, around the Forsyth Range, Maneroo Creek and Longreach, with a note that it comes from the Upper Barcoo. Wikitongues lists it as Pama--Nyungan with an EGIDS status of 10 (extinct). My searches found no current speaker numbers and no teaching materials. Very little exists, and the sources below are only reference records.

\section{Primary Educational Resources}
No primers, textbooks or school programmes for Yiningayi turned up in the search results.

\section{Linguistic and Reference Materials}
\begin{itemize}
\item AIATSIS Austlang record L41, Iningai: \url{https://collection.aiatsis.gov.au/austlang/language/L41}
\item The search summary says Gavan Breen's 1990 study (not seen directly) concluded that the language is closest to Bidjara, Dharawala and Wadjabangayi, and that the data quality is poor. The same summary says Breen compiled a wordlist from Alice River materials. I did not see the study itself.
\item A journal article, ``A partial vocabulary of the Yinikurtira'', appeared in the search results. I could not confirm that it concerns this language, so treat it as a lead only: \url{https://researchportal.murdoch.edu.au/esploro/outputs/journalArticle/A-partial-vocabulary-of-the-Yinikurtira/991005544254107891}
\item University of Sydney archive records surfaced in the results, for example \url{https://archives-search.sydney.edu.au/nodes/view/17507}. I did not verify their content.
\item ASJP wordlist entry (Nyigina): \url{https://asjp.clld.org/languages/NYIGINA}. This is probably a different language, Nyikina, so do not use it.
\item IANA registration for the code \texttt{ygi}: \url{https://www.iana.org/assignments/lang-subtags-templates/ygi.txt}
\end{itemize}

\section{Religious and Community Resources}
None found.

\section{Digital and Audio Resources}
Wikitongues has a language page: \url{https://wikitongues.org/languages/ygi/}. I did not confirm that it holds any recordings.

\section{Curriculum Development Context}
The language is extinct and was poorly documented. A K-12 heritage curriculum would have to be built from archival wordlists, in consultation with the descendant community and AIATSIS. It could rely on documentation of neighbouring related languages only with great caution. No native-language classroom materials exist in what I found.

\end{document}